% !TeX root = ../tesis.tex

\chapter{Anteproyecto}
% \addcontentsline{toc}{chapter}{\protect\numberline{}Introduction}	  		% Comment if you don't want the introduction to appear on the table of content. It will not have a number
\label{chapter:intro}

\section{Planteamiento del problema}
\label{section:planteamiento}

El Juego de la Vida es uno de los autómatas celulares más conocidos, ya que permite observar la evolución y variabilidad de un patrón inicial a través de un conjunto de reglas simples. Esto ha sido de interés para muchos investigadores debido a que se ha utilizado como modelo para estudiar fenómenos complejos en diversas áreas, como la biología, la física, la economía y demás campos \cite{argun2021simulation}\cite{adamatzky2010game}. No obstante, la simulación de este juego llega a ser computacionalmente costosa cuando se trabaja con un gran rango de celdas y generaciones, provocando que sea intratable para una sola computadora común. Es este sentido, se requeriría de costosas computadoras especializadas para poder realizar este trabajo extremo en paralelo.

Por lo tanto, se propone la realización de un clúster de computadoras que permita ejecutar el Juego de la Vida en paralelo, aumentando así el poder de procesamiento a comparación de si se utiliza una sola computadora. Esto con la peculiaridad de utilizar computadoras antiguas, permitiendo que sea asequible por su bajo costo.